\section*{Capítulo \ref{ch:b3:lp} --- Os espaços $L^p$}

\begin{solution}{exo:b3:lp:1}
(a) $p = q = 2$ no~\cref{thm:b3:lp:holder}: $\abs{\int f\bar
g\,\dd\mu} \leq \int\abs{fg} \leq \norm f_2\norm g_2$ ---
Cauchy--Schwarz, com igualdade se, e somente se, $\abs f, \abs g$ são
proporcionais e as fases estão alinhadas.

(b) Num espaço de probabilidade, para $p \leq q$: aplique Jensen
(\cref{exo:b3:lp:10}) com a convexa $\Phi(t) =
\abs t^{q/p}$ à função $\abs f^p$:
$\bigl(\int\abs f^p\bigr)^{q/p} \leq \int\abs f^q$, isto é,
$\norm f_p \leq \norm f_q$.

(c) $\int\abs{fg} = \norm f_1\norm g_\infty$ se, e somente se, $\abs{g} =
\norm g_\infty$ q.t.p.\ em $\{f \neq 0\}$ (a desigualdade
$\abs{fg} \leq \abs f\norm g_\infty$ deve ser uma igualdade
q.t.p.).
\end{solution}

\begin{solution}{exo:b3:lp:2}
$\int_0^1 x^{-p/2}\dd x < \infty$ se, e somente se, $p < 2$: a primeira está em
$L^p$ para $p \in \intco12$. $\int_1^\infty x^{-p/2}\dd x <
\infty$ se, e somente se, $p > 2$: a segunda para $p \in \intoo2\infty$ (e
$p = \infty$: ela é limitada --- inclua-o). Terceira: para $p <
2$, dominada por $x^{-p/2}$: \omterm{def:b3:lebesgue:l1}{integrável}; para $p = 2$,
substitua $u = -\ln x$: $\int_0^1\frac{\dd x}{x(1 + \abs{\ln
x})^2} = \int_0^\infty\frac{\dd u}{(1 + u)^2} < \infty$; para
$p > 2$ a potência domina: divergente. Logo $p \in \intcc12$.
Quarta: $\int_\R\frac{\dd x}{(1 + \abs x)^p} < \infty$ se, e somente se, $p >
1$; limitada, de modo que também $L^\infty$: $p \in
\intoc1\infty$. Moral: a \omterm{def:b3:lebesgue:l1}{integrabilidade} em $0$ gosta de $p$ pequeno, e em
$\infty$ de $p$ grande; combinando os dois obstáculos, não há inclusão
entre os espaços $L^p(\R)$.
\end{solution}

\begin{solution}{exo:b3:lp:3}
(a) $\norm{f_n}_p^p = \lambda(I_n) \to 0$ (no nível diádico $k$
o comprimento é $2^{-k}$). Mas todo $x$ está num intervalo de
cada nível diádico: $f_n(x) = 1$ uma infinidade de vezes e $= 0$
uma infinidade de vezes (intervalos do mesmo nível que não contêm
$x$): nenhuma convergência em ponto algum.

(b) $f_{n_k} = \mathbf 1_{\intcc0{2^{-k}}}$ (o primeiro intervalo
de cada nível) converge a $0$ em todo $x > 0$: q.t.p.

(c) q.t.p.\ mas não em $L^1$: $n\mathbf 1_{\intoo0{1/n}} \to 0$
q.t.p., de integral $1$. Em $L^1$, mas em nenhum $L^p$ ($p > 1$): $g_n =
\eu^n\,\mathbf 1_{(0,\ \eu^{-n}/n)}$: $\norm{g_n}_1 = \frac1n
\to 0$, ao passo que $\norm{g_n}_p^p = \eu^{(p-1)n}/n \to \infty$ para
todo $p > 1$.
\end{solution}

\begin{solution}{exo:b3:lp:4}
Superior: $\norm f_p \leq \mu(X)^{1/p}\norm f_\infty$
(\cref{prop:b3:lp:inclusions}(a) com $q = \infty$), e
$\mu(X)^{1/p} \to 1$. Inferior: para $\varepsilon > 0$, $A =
\{\abs f > \norm f_\infty - \varepsilon\}$ tem $\mu(A) > 0$
(definição do sup essencial), e
\[
\norm f_p \geq \Bigl(\int_A \abs f^p\Bigr)^{1/p}
\geq (\norm f_\infty - \varepsilon)\,\mu(A)^{1/p}
\xrightarrow[p\to\infty]{} \norm f_\infty - \varepsilon .
\]
\end{solution}

\begin{solution}{exo:b3:lp:5}
(a) Duas vezes: a conversão $\norm{\tau_hg - g}_p \leq
C^{1/p}\norm{\tau_hg - g}_\infty$ (suporte de \omterm{def:b3:measure:measure}{medida} finita)
degenera para $p = \infty$ apenas no fato de que ali falha a \emph{\omterm{ex:b3:lebesgue:gamma}{densidade} de
$\mathcal C_c$} --- essa é a lacuna real: o passo (2)
do~\cref{thm:b3:lp:density} não tem análogo $L^\infty$.

(b) Se $f$ tem um representante uniformemente \omterm{def:b3:topology:continuity}{contínuo} $g$:
$\norm{\tau_hf - f}_\infty = \sup_x\abs{g(x - h) - g(x)} \to
0$. Reciprocamente, suponha $\norm{\tau_hf - f}_\infty \to 0$. As
regularizações $f_\varepsilon = f * \rho_\varepsilon$ são
\omterm{def:b3:topology:continuity}{contínuas}, e
\[
\norm{f_\varepsilon - f}_\infty
\leq \sup_{\norm y \leq \varepsilon}\norm{\tau_yf -
f}_\infty \longrightarrow 0
\]
(a convolução é uma média de transladados). Cada
$f_\varepsilon$ é uniformemente \omterm{def:b3:topology:continuity}{contínua}
($\norm{\tau_hf_\varepsilon - f_\varepsilon}_\infty \leq
\norm{\tau_hf - f}_\infty$, por média), e um limite uniforme
de funções uniformemente \omterm{def:b3:topology:continuity}{contínuas} é uniformemente \omterm{def:b3:topology:continuity}{contínuo}: $f$ coincide q.t.p.\
com uma função uniformemente \omterm{def:b3:topology:continuity}{contínua}.
\end{solution}

\begin{solution}{exo:b3:lp:6}
Por Hölder, para todo $x$ o integrando $y \mapsto f(x -
y)g(y)$ está em $L^1$ com $\abs{(f*g)(x)} \leq \norm f_p\norm
g_q$: definida em toda parte e limitada. \omterm{def:b3:topology:continuity}{Continuidade} uniforme ($p <
\infty$):
\[
\abs{(f*g)(x + h) - (f*g)(x)}
= \Bigl|\int\bigl(\tau_{-h}f - f\bigr)(x - y)\,g(y)\dd y\Bigr|
\leq \norm{\tau_{-h}f - f}_p\,\norm g_q
\xrightarrow[h\to0]{} 0,
\]
uniformemente em $x$ (\cref{thm:b3:lp:density}(3)). Se $p =
\infty$, então $q = 1$: escreva $f * g = g * f$ e repita a mesma
cota com a translação agindo em $g \in L^1$.
\end{solution}

\begin{solution}{exo:b3:lp:7}
Fixe $\chi \in \mathcal C_c^\infty(\intoo ab)$ com $\int\chi =
1$ e ponha $c = \int f\chi$. Seja $\varphi \in \mathcal
C_c^\infty$ arbitrária e $\psi = \varphi -
\bigl(\int\varphi\bigr)\chi$: então $\int\psi = 0$, de modo que $\Phi(x)
= \int_a^x\psi$ define $\Phi \in \mathcal C_c^\infty(\intoo
ab)$ (ela se anula perto dos dois extremos: perto de $a$ trivialmente, perto de $b$
porque a integral total é $0$) com $\Phi' = \psi$. A
hipótese dá $\int f\psi = \int f\Phi' = 0$, logo
\[
\int f\varphi = \Bigl(\int\varphi\Bigr)\int f\chi = \int
c\,\varphi
\quad\text{para todo } \varphi:
\qquad \int(f - c)\varphi = 0 .
\]
Pelo~\cref{cor:b3:lp:fundlemma} (localizado em $\intoo ab$), $f =
c$ q.t.p.
\end{solution}

\begin{solution}{exo:b3:lp:8}
Seja $3\delta < d(K, \R^d\setminus U)$ (positivo:
\cref{exo:b3:topology:6}(b); se $U = \R^d$, qualquer $\delta$
serve), $K_\delta = \{x : d(x, K) \leq \delta\}$, e
\[
\varphi = \mathbf 1_{K_\delta} * \rho_{\delta/2} .
\]
Então $\varphi \in \mathcal C^\infty$
(\cref{thm:b3:lp:regularization}(1); a indicadora é $L^1$),
$0 \leq \varphi \leq 1$ ($\int\rho = 1$), $\varphi = 1$ em $K$
(para $x \in K$, $\bar B(x, \delta/2) \subseteq K_\delta$, de modo que
a convolução integra $\rho$ por inteiro) e
$\operatorname{supp}\varphi \subseteq K_{3\delta/2} \subseteq
U$: com suporte \omterm{def:b3:topology:compact}{compacto} ($K_\delta$ é limitado). Partição da
unidade: dado $K \subseteq U_1\cup\dots\cup U_m$, escolha (por
\omterm{def:b3:topology:compact}{compacidade}) \omterm{def:b3:topology:compact}{compactos} $K_i \subseteq U_i$ com $K \subseteq
\bigcup \mathring K_i$, tome $\varphi_i$ como acima para $(K_i,
U_i)$ e ponha $\psi_i = \varphi_i\prod_{j <i}(1 -
\varphi_j)$: cada $\psi_i \in \mathcal C_c^\infty(U_i)$, e
$\sum_i\psi_i = 1 - \prod_i(1 - \varphi_i) = 1$ em $K$.
\end{solution}

\begin{solution}{exo:b3:lp:9}
(a) O expoente de interpolação para $(p_0, q_0) = (1, \infty)$
em $r = p$ é $\alpha = \frac1p$:
a~\cref{prop:b3:lp:inclusions}(c) dá $\norm f_p \leq \norm
f_1^{1/p}\norm f_\infty^{1 - 1/p}$.

(b) Tomando logaritmos na~\cref{prop:b3:lp:inclusions}(c): $\ln\norm f_r \leq
\alpha\ln\norm f_p + (1-\alpha)\ln\norm f_q$, em que $\frac1r$
é a mesma combinação convexa de $\frac1p, \frac1q$: $\frac1p
\mapsto \ln\norm f_p$ é convexa. Exemplo com pertencimento a $L^p$
exatamente em $\intoo{p_0}{p_1}$:
\[
f(x) = x^{-1/p_1}\,\mathbf 1_{\intoo01}(x) +
x^{-1/p_0}\,\mathbf 1_{\intco1\infty}(x):
\]
o primeiro termo é $L^p$ se, e somente se, $p < p_1$, e o segundo se, e somente se, $p >
p_0$.
\end{solution}

\begin{solution}{exo:b3:lp:10}
Seja $m = \int f\,\dd\mu \in \R$. A convexidade fornece uma reta de
suporte em $m$: existe $s$ com $\Phi(t) \geq \Phi(m) + s(t -
m)$ para todo $t$ (tome $s$ entre as derivadas laterais,
que existem para funções convexas). Substitua $t = f(x)$ e
integre contra a probabilidade $\mu$:
\[
\int\Phi\circ f\,\dd\mu \geq \Phi(m) + s\Bigl(\int f - m\Bigr)
= \Phi\Bigl(\int f\,\dd\mu\Bigr)
\]
(\omterm{def:b3:lebesgue:measurable}{mensurabilidade}: $\Phi$ é \omterm{def:b3:topology:continuity}{contínua}; a \omterm{def:b3:lebesgue:l1}{integrabilidade} da
parte negativa de $\Phi\circ f$ é garantida pela reta de
suporte). MA--MG: num conjunto finito com pesos $w_i$, tome
$\Phi = \exp$ e $f = \sum(\ln a_i)\mathbf 1_i$:
$\exp\bigl(\sum w_i\ln a_i\bigr) \leq \sum w_ia_i$, isto é,
$\prod a_i^{w_i} \leq \sum w_ia_i$. A monotonicidade das normas é
o~\cref{exo:b3:lp:1}(b).
\end{solution}

\begin{solution}{exo:b3:lp:11}
(a) Suponha primeiro $p, q, r < \infty$ e $f, g \geq 0$
(substitua pelos valores absolutos). Os três expoentes $r$,
$\alpha = \frac{pr}{r-p}$, $\beta = \frac{qr}{r-q}$
satisfazem $\frac1r + \frac1\alpha + \frac1\beta = \frac1r +
\frac1p - \frac1r + \frac1q - \frac1r = 1$ (a relação de
escala). Reparta, para $x$ fixo,
\[
f(y)g(x-y) = \bigl[f(y)^pg(x-y)^q\bigr]^{1/r}\cdot
f(y)^{1 - p/r}\cdot g(x-y)^{1 - q/r},
\]
e Hölder com os três expoentes dá
\[
(f*g)(x) \leq \Bigl(\int f^pg(x-\cdot)^q\Bigr)^{1/r}
\norm f_p^{\,p/\alpha\cdot\alpha/p}\cdots
\]
mais precisamente: o segundo fator é
$\bigl(\int f^{(1-p/r)\alpha}\bigr)^{1/\alpha} =
\norm f_p^{p(1/p - 1/r)}$, pois $(1 - \frac pr)\alpha = p$,
e do mesmo modo o terceiro é $\norm g_q^{q(1/q - 1/r)}$. Eleve
à $r$-ésima potência e integre em $x$ (Tonelli no
primeiro fator):
\[
\norm{f*g}_r^r \leq \norm f_p^p\,\norm g_q^q\cdot
\norm f_p^{\,rp(1/p - 1/r)}\,\norm g_q^{\,rq(1/q - 1/r)}
= \norm f_p^r\,\norm g_q^r .
\]
Os casos extremos ($r = \infty$ ou um expoente igual a
sua cota) são Hölder puro ou estimativas diretas.

(b) $r = \infty$ força $q = p'$: $\abs{f*g(x)} \leq \norm
f_p\norm g_{p'}$ --- Hölder após translação e reflexão.
$q = 1$ dá $r = p$: $\norm{f*g}_p \leq \norm
g_1\norm f_p$, o burro de carga da regularização
(o motor do~\cref{thm:b3:lp:regularization}). $p = q = 1$
dá $r = 1$: a álgebra de convolução
(\cref{thm:b3:product:convolution}).

(c) Substitua $f, g$ por $f_\lambda = f(\lambda\cdot)$,
$g_\lambda = g(\lambda\cdot)$: então $f_\lambda*g_\lambda =
\lambda^{-d}(f*g)(\lambda\cdot)$ e, comparando normas,
\[
\text{ME} \sim \lambda^{-d - d/r}, \qquad
\text{MD} \sim \lambda^{-d/p - d/q} :
\]
uma desigualdade válida para todos $f, g$ força os dois expoentes de
escala a coincidir, isto é, $1 + \frac1r = \frac1p + \frac1q$
exatamente. Qualquer outra combinação morre em $\lambda \to 0$ ou em
$\infty$: a relação de Young não é uma conveniência, mas uma
lei de escala.
\end{solution}

\begin{solution}{exo:b3:lp:12}
(a) Normalize $\norm f_p = \norm g_q = 1$. A demonstração de
Hölder integra a desigualdade de Young $\abs{fg} \leq
\frac{\abs f^p}p + \frac{\abs g^q}q$; a igualdade das
integrais força a igualdade q.t.p.\ em Young, o que (convexidade
estrita de $\exp$; igualdade se, e somente se, $a^p = b^q$) significa $\abs
f^p = \abs g^q$ q.t.p. Desfazendo a normalização:
$\abs f^p\norm g_q^q = \abs g^q\norm f_p^p$ q.t.p.\ ---
proporcionalidade.

(b) Minkowski são dois Hölder aplicados a $\abs{f + g}^{p-1}
\abs f$ e $\abs{f+g}^{p-1}\abs g$; a igualdade força a
proporcionalidade de (a) em ambos: $\abs f^p$ e $\abs g^p$ cada um
proporcional a $\abs{f+g}^{(p-1)q} = \abs{f+g}^p$, de modo que
$\abs g = t\abs f$ q.t.p.\ para uma constante $t \geq 0$; e a
desigualdade triangular pontual inicial $\abs{f + g} \leq
\abs f + \abs g$ também tem de ser uma igualdade q.t.p., o que, para
valores complexos, significa que $f$ e $g$ têm q.t.p.\ o mesmo
argumento onde ambas são não nulas. Combinando: $g = tf$ q.t.p.,
$t > 0$ (ambas não nulas).

(c) $p = 1$: a igualdade em $\int\abs{f + g} = \int\abs f +
\int\abs g$ vale sempre que $f, g$ têm o mesmo padrão de
sinais (mesmo argumento q.t.p.) --- sem necessidade de proporcionalidade:
$f = \mathbf 1_{\intcc01}$ e $g = \mathbf 1_{\intcc02}$
servem. $p = \infty$: $\norm{f+g}_\infty = \norm f_\infty +
\norm g_\infty$ assim que as duas funções atingem seus picos de modo
compatível num ponto comum (ou ao longo de uma sequência comum):
$f = g\,$ perto de um ponto basta, qualquer que seja o comportamento
nos demais. A convexidade estrita das bolas $L^p$ para $1 <
p < \infty$ --- e sua falha nos extremos --- é
exatamente o que esses casos de igualdade testemunham.
\end{solution}

\begin{solution}{pb:b3:lp:1}
\textbf{1.} $f$ tem suporte em algum $[\alpha, \beta]
\subseteq \intoo0\infty$, de modo que $F = 0$ em $[0, \alpha]$ e $F
\equiv F(\beta)$ em $[\beta, \infty)$: $Hf$ se anula perto de $0$
e é $O(1/x)$ no infinito; $\int_1^\infty x^{-p}\dd x <
\infty$ para $p > 1$, e $Hf$ é \omterm{def:b3:topology:continuity}{contínua}: $Hf \in L^p$.

\textbf{2.} $\bigl(x^{1-p}F(x)^p\bigr)' = (1-p)x^{-p}F^p +
p\,x^{1-p}F^{p-1}f$. Os dois valores de bordo se anulam: em $0$
porque $F = 0$ perto de $0$; em $\infty$ porque $x^{1-p}F^p \leq
F(\beta)^p x^{1-p} \to 0$ ($p > 1$). Integrando a identidade
em $\intoo0\infty$:
\[
0 = (1 - p)\int_0^\infty\Bigl(\frac Fx\Bigr)^p\dd x
+ p\int_0^\infty\Bigl(\frac Fx\Bigr)^{p-1}f(x)\,\dd x,
\]
que é a relação em destaque.

\textbf{3.} Hölder com expoentes $q = \frac p{p-1}$ e $p$:
\[
\int\Bigl(\frac Fx\Bigr)^{p-1}f
\leq \Bigl(\int\Bigl(\frac
Fx\Bigr)^{p}\Bigr)^{1 - 1/p}\,\norm f_p ,
\]
de modo que $\norm{Hf}_p^p \leq \frac p{p-1}\norm{Hf}_p^{p-1}\norm
f_p$; divida (finito pela questão 1 e, se for $0$, nada há
a demonstrar).

\textbf{4.} Sejam $f \in L^p$, $f \geq 0$. Escolha $\varphi_k \in
\mathcal C_c(\intoo0\infty)$ com $\varphi_k \to f$ em $L^p$
(\cref{thm:b3:lp:density}(2), intersectado com a semirreta
aberta --- aproxime $f\mathbf 1_{[1/k, k]}$ e
diagonalize), e substitua $\varphi_k$ por $\abs{\varphi_k}$
(ainda \omterm{def:b3:topology:continuity}{contínua} e mais próxima de $f \geq 0$). Para todo
$x > 0$ fixo, Hölder em $\intoo0x$ dá
\[
\abs{H\varphi_k(x) - Hf(x)} \leq
\frac1x\,x^{1 - 1/p}\,\norm{\varphi_k - f}_p \to 0 :
\]
$H\varphi_k \to Hf$ pontualmente. Fatou e a questão 3:
\[
\int (Hf)^p \leq \liminf_k\int(H\varphi_k)^p
\leq \Bigl(\frac p{p-1}\Bigr)^p\liminf_k\norm{\varphi_k}_p^p
= \Bigl(\frac p{p-1}\Bigr)^p\norm f_p^p .
\]
Para $f$ com sinal ou complexa: $\abs{Hf} \leq H\abs f$ pontualmente,
e o caso não negativo se aplica a $\abs f$.

\textbf{5.} $\norm{f_A}_p^p = \int_1^A t^{-1}\dd t = \ln A$.
Para $1 \leq x \leq A$:
\[
(Hf_A)(x) = \frac1x\int_1^x t^{-1/p}\dd t
= \frac{x^{1 - 1/p} - 1}{(1 - \tfrac1p)\,x}
= \frac p{p-1}\,x^{-1/p}\bigl(1 - x^{-(1 - 1/p)}\bigr).
\]

\textbf{6.} Fixe $\varepsilon > 0$ e $X_0$ com $(1 -
x^{-(1-1/p)})^p \geq 1 - \varepsilon$ para $x \geq X_0$. Então
\[
\norm{Hf_A}_p^p \geq \int_{X_0}^A\Bigl(\frac
p{p-1}\Bigr)^p\frac{1 - \varepsilon}{x}\,\dd x
= \Bigl(\frac p{p-1}\Bigr)^p(1 - \varepsilon)\,(\ln A - \ln
X_0),
\]
de modo que $\dfrac{\norm{Hf_A}_p^p}{\norm{f_A}_p^p} \geq \bigl(\frac
p{p-1}\bigr)^p(1 - \varepsilon)\bigl(1 - \frac{\ln X_0}{\ln
A}\bigr) \to \bigl(\frac p{p-1}\bigr)^p(1 - \varepsilon)$ quando
$A \to \infty$: a constante não pode ser melhorada.

\textbf{7.} A igualdade na questão 3 força a igualdade em
Hölder: $f^p$ proporcional a $\bigl(\frac Fx\bigr)^{(p-1)q}
= \bigl(\frac Fx\bigr)^p$ q.t.p., isto é, $f = \gamma\,\frac Fx$
q.t.p.\ para algum $\gamma \geq 0$. Como $F(x) = \int_0^xf$ é
absolutamente \omterm{def:b3:topology:continuity}{contínua} com $F' = f$ q.t.p., $F$ resolve $F' =
\gamma F/x$: em qualquer intervalo em que $F > 0$, $(\ln F)' =
\gamma/x$, de modo que $F = c\,x^{\gamma}$ e $f = c\gamma
x^{\gamma - 1}$ ali. Mas nenhuma potência não nula $x^{\gamma-1}$
pertence a $L^p(\intoo0\infty)$ (é preciso $p(\gamma - 1) < -1$
em $\infty$ e $> -1$ em $0$: incompatível), e $F$ não pode
se anular identicamente a menos que $f = 0$. Logo a igualdade exige $f =
0$.

\textbf{8.} Dada $(a_n)$ não crescente $\geq 0$, defina a
função escada $g(t) = a_{\lceil t\rceil}$ em
$\intoo0{+\infty}$: não crescente, com $\int_0^\infty g^p =
\sum_na_n^p$ e $\int_0^ng = a_1 + \dots + a_n$, de modo que
$(Hg)(n) = \frac{a_1 + \dots + a_n}n$. A média de uma função
não crescente é não crescente, de modo que $Hg$ o é, e
\[
\sum_{n\geq1}\Bigl(\frac{a_1{+}\dots{+}a_n}n\Bigr)^p
= \sum_{n\geq1}(Hg)(n)^p
\leq \sum_{n\geq1}\int_{n-1}^n (Hg)(t)^p\,\dd t
= \norm{Hg}_p^p
\leq \Bigl(\frac p{p-1}\Bigr)^p\sum_n a_n^p
\]
pela Parte I. Para uma sequência não negativa geral, seja $(a_n^*)$
seu rearranjo não crescente (possível quando $a_n \to 0$,
o que podemos supor --- caso contrário, ambos os membros são infinitos):
o membro direito não muda, e cada soma parcial $a_1 + \dots
+ a_n$ é no máximo $a_1^* + \dots + a_n^*$ (os $n$ maiores
termos): o membro esquerdo apenas cresce. Logo a desigualdade vale para toda
$(a_n)$.

\textbf{9.} Se $(a_n) \in \ell^p$, a sequência das médias de Cesàro
está em $\ell^p$ com norma $\leq \frac p{p-1}\norm
a_p$. Exemplo ($p = 2$): $a_n = \frac1{\sqrt n\,\ln n}$ ($n
\geq 2$): $\sum a_n^2 = \sum\frac1{n\ln^2n} < \infty$, e, no entanto,
$\sum a_n = \infty$ (teste da integral); Hardy ainda garante
$\sum_n\bigl(\frac{a_1 + \dots + a_n}n\bigr)^2 < \infty$.

\textbf{10.} Para $f = \mathbf 1_{\intcc01}$: $Hf(x) = 1$ em
$\intoc01$ e $= \frac1x$ para $x \geq 1$: $\int_0^\infty Hf =
1 + \int_1^\infty\frac{\dd x}x = \infty$, ao passo que $\norm f_1 =
1$. A demonstração desaba em dois pontos: a constante
$\frac p{p-1}$ explode quando $p \to 1$, e o termo de bordo
$x^{1-p}F^p$ deixa de se anular no infinito para $p = 1$.
A desigualdade de Hardy é um fenômeno honestamente de $p > 1$.

\textbf{11.} Tome o intervalo mais longo $B_{i_1}$; descarte
todo intervalo que o encontre; tome o mais longo sobrevivente
$B_{i_2}$; itere (há finitos intervalos). Os escolhidos
são disjuntos por construção, e todo $B$ descartado encontrou um
intervalo escolhido ao menos tão longo: um intervalo que encontra um
de comprimento maior ou igual está contido em seu triplo, $B \subseteq
3B_{i_l}$.

\textbf{12.} $\{Mf > t\}$ é aberto: cada média $x \mapsto
\frac1{2r}\int_{x-r}^{x+r}\abs f$ é \omterm{def:b3:topology:continuity}{contínua} (convergência
dominada em $x$), e um supremo de funções \omterm{def:b3:topology:continuity}{contínuas}
é semicontínuo inferiormente. Cada $x$ nele possui $B_x =
\intoo{x-r_x}{x+r_x}$ com $\int_{B_x}\abs f >
t\,\lambda(B_x)$. Para $K \subseteq \{Mf > t\}$ \omterm{def:b3:topology:compact}{compacto}:
finitos $B_x$ cobrem $K$, Vitali (questão 11) extrai
$B_1', \dots, B_k'$ disjuntos com $K \subseteq
\bigcup_l3B_l'$, de modo que
\[
\lambda(K) \leq 3\sum_l\lambda(B_l') <
\frac3t\sum_l\int_{B_l'}\abs f \leq \frac3t\,\norm f_1
\]
por disjunção; a regularidade \omterm{def:b3:topology:interior}{interior}
(\cref{thm:b3:measure:regularity}) conclui.

\textbf{13.} Para $x > 1$: com $r \in \intcc{x-1}x$, a
média é $\frac{r - x + 1}{2r}$, crescente em $r$; para $r
\geq x$ ela é $\frac1{2r}$, decrescente: o supremo é
$\frac1{2x}$, atingido em $r = x$. Logo $M\mathbf 1_{\intcc01}
\notin L^1$. Em geral, se $\int_I\abs f = c > 0$ num
intervalo limitado $I \subseteq \intcc{-C}C$, então $Mf(x) \geq
\frac{c}{2(\abs x + C)}$ para todo $x$: nunca \omterm{def:b3:lebesgue:l1}{integrável},
a menos que $f = 0$ q.t.p.

\textbf{14.} Com $f_t = f\,\mathbf 1_{\abs f > t/2}$:
$M(f - f_t) \leq \frac t2$, de modo que $\{Mf > t\} \subseteq \{Mf_t
> \frac t2\}$ e a questão 12 dá $\lambda(Mf > t) \leq
\frac6t\int_{\abs f > t/2}\abs f$. Fórmula das camadas
(\cref{prop:b3:product:layercake}) e Tonelli:
\[
\norm{Mf}_p^p = p\int_0^\infty t^{p-1}\lambda(Mf > t)\dd t
\leq 6p\int\abs f\int_0^{2\abs f}t^{p-2}\,\dd t\,\dd\lambda
= \frac{6p\,2^{p-1}}{p-1}\,\norm f_p^p .
\]

\textbf{15.} Escreva $A_rf(x) = \frac1{2r}\int_{x-r}^{x+r}f$.
Para $g$ \omterm{def:b3:topology:continuity}{contínua}: $A_rg(x) \to g(x)$ em toda parte. Dado
$\varepsilon > 0$, reparta $f = g + h$ com $g$ \omterm{def:b3:topology:continuity}{contínua} de
suporte \omterm{def:b3:topology:compact}{compacto} e $\norm h_1 < \varepsilon$
(\cref{thm:b3:lp:density}); então
\[
\limsup_{r\to0}\,\abs{A_rf(x) - f(x)} \leq Mh(x) +
\abs{h(x)},
\]
de modo que $\Omega_\delta = \{\limsup_r\abs{A_rf - f} > \delta\}
\subseteq \{Mh > \tfrac\delta2\} \cup \{\abs h >
\tfrac\delta2\}$ tem \omterm{def:b3:measure:measure}{medida} $\leq \frac{6\varepsilon}\delta
+ \frac{2\varepsilon}\delta$ (questão 12; Markov).
$\varepsilon$ arbitrário: $\lambda(\Omega_\delta) = 0$; união
em $\delta = \frac1k$: $A_rf \to f$ q.t.p.

\textbf{16.} (a) Para cada $q \in \Q$, a questão 15 aplicada a
$\abs{f - q}$ dá $A_r\abs{f - q}(x) \to \abs{f(x) - q}$
q.t.p.; na interseção desses conjuntos de \omterm{def:b3:measure:measure}{medida} total, escolha
$q$ com $\abs{f(x) - q} < \eta$:
$\limsup_rA_r\abs{f - f(x)}(x) \leq 2\eta$ para todo $\eta$:
quase todo $x$ é ponto de Lebesgue. (b) Num ponto de Lebesgue,
\[
\Bigl|\frac{F(x + h) - F(x)}h - f(x)\Bigr|
= \Bigl|\frac1h\int_x^{x+h}\bigl(f - f(x)\bigr)\Bigr|
\leq 2\,A_{\abs h}\abs{f - f(x)}(x) \to 0 :
\]
$F' = f$ q.t.p.\ --- primitivas de funções $L^1$ de fato
derivam de volta; a escada (\cref{pb:b3:measure:1})
é o contraexemplo apenas para a direção \emph{recíproca}.

\textbf{17.} Aplique a questão 15 a $\mathbf 1_{A\cap[-n,n]}$
e faça $n$ crescer: para q.t.p.\ $x \in A$, a \omterm{ex:b3:lebesgue:gamma}{densidade}
$\frac{\lambda(A\cap\intcc{x-r}{x+r})}{2r} \to 1$. Steinhaus:
em torno de um ponto de \omterm{ex:b3:lebesgue:gamma}{densidade}, tome $r$ com \omterm{ex:b3:lebesgue:gamma}{densidade} $> \frac34$;
para $\abs t < \frac r2$, $A$ e $A + t$ preenchem, cada um, mais de
$\frac32r$ de um intervalo de comprimento $\leq \frac52r$, e portanto
se intersectam: $\intoo{-r/2}{r/2} \subseteq A - A$.

\textbf{18.} Seja $G(x) = x^{\alpha+1-p}F(x)^p$: $G$ se anula
em $0$ ($F$ se anula perto de $0$) e em $\infty$ ($F$ limitada,
$\alpha + 1 - p < 0$), de modo que $\int_0^\infty G' = 0$ com
\[
G' = (\alpha + 1 - p)\,x^{\alpha-p}F^p +
p\,x^{\alpha+1-p}F^{p-1}f
= (\alpha + 1 - p)\Bigl(\frac Fx\Bigr)^px^\alpha +
p\Bigl(\frac Fx\Bigr)^{p-1}f\,x^\alpha .
\]
Logo $\int(F/x)^px^\alpha = \frac{p}{p-1-\alpha}
\int(F/x)^{p-1}f\,x^\alpha$; Hölder para a \omterm{def:b3:measure:measure}{medida}
$x^\alpha\dd x$ (expoentes $\frac p{p-1}$ e $p$) conclui
como na questão 3. Caso limite $\alpha = p - 1$: com $f(t) =
\frac1t\mathbf 1_{\intcc1A}$, o membro direito é $\ln A$,
ao passo que o esquerdo contém $\int_1^A\frac{(\ln x)^p}x\dd x =
\frac{(\ln A)^{p+1}}{p+1}$: nenhuma constante sobrevive a $A \to
\infty$.

\textbf{19.} Tonelli em $\{0 < t < x\}$:
\[
\langle Hf, g\rangle =
\int_0^\infty\frac{g(x)}x\int_0^xf(t)\,\dd t\,\dd x
= \int_0^\infty f(t)\int_t^\infty\frac{g(x)}x\,\dd x\,\dd t
= \langle f, H^*g\rangle .
\]
Dualidade: $\norm{H^*f}_p = \sup\{\langle f, Hg\rangle : g
\geq 0, \norm g_q \leq 1\} \leq \norm f_p\cdot
\sup\norm{Hg}_q \leq \frac q{q-1}\norm f_p = p\,\norm f_p$,
sendo Hardy invocado em $L^q$, cuja constante $\frac q{q-1}$
é igual a $p$.

\textbf{20.} Reparta ao longo da diagonal (nula). Em $\{y \leq
x\}$:
\[
\iint_{y\leq x}\frac{f(x)g(y)}{x}\,\dd y\,\dd x
= \int f(x)\,(Hg)(x)\,\dd x \leq \norm f_p\,\norm{Hg}_q
\leq p\,\norm f_p\norm g_q,
\]
pois Hardy em $L^q$ carrega a constante $\frac q{q-1} =
p$. Simetricamente, $\iint_{y>x} = \int g\,(Hf) \leq
\norm g_q\norm{Hf}_p \leq q\,\norm f_p\norm g_q$ ($\frac
p{p-1} = q$). Total: $(p + q)\,\norm f_p\norm g_q$.

\textbf{21.} Para $a_n = n^{-1/p}$, $n \leq N$: o membro
direito é $\bigl(\frac p{p-1}\bigr)^p\sum_{n\leq N}\frac1n =
\bigl(\frac p{p-1}\bigr)^p\ln N + O(1)$. No esquerdo, para $n
\leq N$: $a_1 + \dots + a_n \geq \int_1^{n+1}t^{-1/p}\dd t =
\frac{p}{p-1}\bigl((n+1)^{1-1/p} - 1\bigr)$, de modo que a $n$-ésima
média de Cesàro é $\geq \frac p{p-1}n^{-1/p}(1 -
o(1))$ uniformemente para $n$ em qualquer faixa $n \geq n_0(\eta)$;
elevando à $p$ e somando, o membro esquerdo é $\geq
\bigl(\frac p{p-1}\bigr)^p(1 - \eta)\ln N + O_\eta(1)$.
Dividindo e fazendo $N \to \infty$ e depois $\eta \to 0$: nenhuma
constante menor que $\bigl(\frac p{p-1}\bigr)^p$ pode servir.

\textbf{22.} $H$ limitado em $L^p$: médias cumulativas não
inflam as normas $p$ (constante $\frac p{p-1}$); Cesàro em
$\ell^p$: o mesmo, discretizado; $M$ limitado em $L^p$: nem
mesmo a melhor média local escapa ao controle (constante
$O(\frac1{p-1})$). Os três falham em $p = 1$, e por uma
única razão: fazer a média de uma unidade de massa concentrada produz uma
cauda $\frac1x$ (questões 10 e 13), e $\frac1x$ pertence
a todo $L^p$ perto do infinito, exceto $L^1$. A suavização espalha
a massa exatamente até a fronteira harmônica da \omterm{def:b3:lebesgue:l1}{integrabilidade}.

\textbf{23.} Ponha $b_n = a_n^{1/p}$, de modo que $\sum b_n^p = \sum
a_n$. A questão 8 dá
\[
\sum_{n\geq1}\Bigl(\frac{b_1 + \dots + b_n}n\Bigr)^{p}
\leq \Bigl(\frac p{p-1}\Bigr)^{p}\sum_{n\geq1}a_n,
\]
e a MA--MG minora cada parcela:
\[
\Bigl(\frac{b_1 + \dots + b_n}n\Bigr)^{p}
\geq \bigl(b_1\cdots b_n\bigr)^{p/n}
= \bigl(a_1\cdots a_n\bigr)^{1/n}.
\]
Logo $\sum(a_1\cdots a_n)^{1/n} \leq \bigl(\frac
p{p-1}\bigr)^p\sum a_n$ para \emph{todo} $p > 1$. Com $m =
p - 1$,
\[
\Bigl(\frac p{p-1}\Bigr)^{p} =
\Bigl(1 + \frac1m\Bigr)^{m+1},
\]
que decresce para $\eu$ quando $m \to \infty$ (a clássica
sequência monótona superior para $\eu$). Tomando o ínfimo em
$p$, obtém-se a desigualdade de Carleman com constante $\eu$.

\textbf{24.} Para $a_n = \frac1n$, $(a_1\cdots a_n)^{1/n} =
(n!)^{-1/n}$. O enquadramento do~\cref{pb:b3:product:1} dá
$\sqrt{2\pi n}\,(n/\eu)^n \leq n! \leq \sqrt{2\pi
n}\,(n/\eu)^n\eu^{1/(12n)}$, de modo que
\[
(n!)^{1/n} = \frac n\eu\,(2\pi n)^{1/(2n)}
\eu^{O(1/n^2)}
= \frac n\eu\Bigl(1 +
O\Bigl(\frac{\ln n}{n}\Bigr)\Bigr),
\]
pois $(2\pi n)^{1/(2n)} = \exp\bigl(\frac{\ln(2\pi
n)}{2n}\bigr)$. Invertendo, $(n!)^{-1/n} = \frac\eu n(1 +
O(\frac{\ln n}n))$, e somando em $n \leq N$:
\[
\sum_{n\leq N}(n!)^{-1/n}
= \eu\ln N + O(1),
\qquad
\eu\sum_{n\leq N}\frac1n = \eu\ln N + O(1)
\]
(a série de erros $\sum\frac{\ln n}{n^2}$ converge). Uma
desigualdade de Carleman com constante $c$ forçaria $\eu\ln N
+ O(1) \leq c\,(\ln N + O(1))$, logo $c \geq \eu$ ao
dividir por $\ln N$. As sequências otimizantes se alinham:
a constante de Hardy é aproximada por $n^{-1/p}$ (questão 21),
a de Carleman por $n^{-1}$ --- em cada caso, a sequência do tipo
harmônico que fica logo fora do espaço sobre o qual se faz a média.

\textbf{25.} Para $f = \mathbf 1_{\intcc01}$: $Hf(x) = 1$ em
$\intoc01$ e $Hf(x) = \frac1x$ para $x > 1$, de modo que
$\norm{Hf}_2^2 = 1 + \int_1^\infty x^{-2}\dd x = 2$ e a
razão é $\sqrt2 \approx 1.414$, cerca de $71\%$ da cota
ótima. Para $f_A$ ($p = 2$): $F(x) = 2(\sqrt x - 1)$ em
$\intcc1A$, de modo que, nessa faixa, $Hf_A = 2x^{-1/2} - 2x^{-1}$
e, para $x > A$, $Hf_A(x) = 2(\sqrt A - 1)/x$. Elevando ao quadrado
e integrando,
\begin{align*}
\int_1^A\bigl(2x^{-1/2} - 2x^{-1}\bigr)^2\dd x
&= 4\ln A - 12 + \frac{16}{\sqrt A} - \frac4A,\\
\int_A^{\infty}\frac{4(\sqrt A - 1)^2}{x^2}\,\dd x
&= 4 - \frac{8}{\sqrt A} + \frac4A,
\end{align*}
donde $\norm{Hf_A}_2^2 = 4\ln A - 8 + 8A^{-1/2}$; dividir
por $\norm{f_A}_2^2 = \ln A$ dá a identidade enunciada. Em $A
= \eu^{10}$: $4 - \frac{8(1 - \eu^{-5})}{10} = 3.2054$, de modo que
a razão é $\sqrt{3.2054} \approx 1.790 < 2$. O defeito
$4 - \norm{Hf_A}_2^2/\norm{f_A}_2^2 \sim 8/\ln A$ decai
apenas logaritmicamente: para atingir a razão $1.99$ seria preciso
$\ln A \approx 200$, isto é, $A \approx 10^{87}$. Constantes
ótimas são teoremas, não experimentos.
\end{solution}
